\documentclass[10pt,a4paper]{amsart}
\usepackage{amsmath,amssymb,amsthm,graphicx}
% Retained preamble for the Stacks Project chapter snapshot in this directory.
% The chapter's own preamble.tex is a copy of the shared Stacks Project preamble; the block
% below is sliced MECHANICALLY, line for line, from it. Nothing was added, removed or
% redefined; the wrapper class, the xr-hyper/hyperref wiring and the externaldocument list
% are omitted because the assembler supplies the document class.
%
% lines taken from preamble.tex: 3, 5, 9, 11, 12, 156, 157, 158, ... (51 lines total)
%
% package coverage restored from the chapter's own preamble (2026-10-13): xy, xy-2cell

\usepackage[all]{xy}
\xyoption{2cell}
\UseAllTwocells

\newenvironment{reference}{\comment}{\endcomment}
\newenvironment{slogan}{\comment}{\endcomment}
\newenvironment{history}{\comment}{\endcomment}
\theoremstyle{plain}
\newtheorem{theorem}[subsection]{Theorem}
\newtheorem{proposition}[subsection]{Proposition}
\newtheorem{lemma}[subsection]{Lemma}
\theoremstyle{definition}
\newtheorem{definition}[subsection]{Definition}
\newtheorem{example}[subsection]{Example}
\newtheorem{exercise}[subsection]{Exercise}
\newtheorem{situation}[subsection]{Situation}
\theoremstyle{remark}
\newtheorem{remark}[subsection]{Remark}
\newtheorem{remarks}[subsection]{Remarks}
\numberwithin{equation}{subsection}
\def\lim{\mathop{\mathrm{lim}}\nolimits}
\def\colim{\mathop{\mathrm{colim}}\nolimits}
\def\Spec{\mathop{\mathrm{Spec}}}
\def\Hom{\mathop{\mathrm{Hom}}\nolimits}
\def\Ext{\mathop{\mathrm{Ext}}\nolimits}
\def\SheafHom{\mathop{\mathcal{H}\!\mathit{om}}\nolimits}
\def\SheafExt{\mathop{\mathcal{E}\!\mathit{xt}}\nolimits}
\def\Sch{\mathit{Sch}}
\def\Mor{\mathop{\mathrm{Mor}}\nolimits}
\def\Ob{\mathop{\mathrm{Ob}}\nolimits}
\def\Sh{\mathop{\mathit{Sh}}\nolimits}
\def\NL{\mathop{N\!L}\nolimits}
\def\CH{\mathop{\mathrm{CH}}\nolimits}
\def\proetale{{pro\text{-}\acute{e}tale}}
\def\etale{{\acute{e}tale}}
\def\QCoh{\mathit{QCoh}}
\def\Ker{\mathop{\mathrm{Ker}}}
\def\Im{\mathop{\mathrm{Im}}}
\def\Coker{\mathop{\mathrm{Coker}}}
\def\Coim{\mathop{\mathrm{Coim}}}
\DeclareMathSymbol{\boxtimes}{\mathbin}{AMSa}{"02}
\def\QCohstack{\mathcal{QC}\!\mathit{oh}}
\def\Cohstack{\mathcal{C}\!\mathit{oh}}
\def\Spacesstack{\mathcal{S}\!\mathit{paces}}
\def\Quotfunctor{\mathrm{Quot}}
\def\Hilbfunctor{\mathrm{Hilb}}
\def\Curvesstack{\mathcal{C}\!\mathit{urves}}
\def\Polarizedstack{\mathcal{P}\!\mathit{olarized}}
\def\Complexesstack{\mathcal{C}\!\mathit{omplexes}}
\def\Pic{\mathop{\mathrm{Pic}}\nolimits}
\def\Picardstack{\mathcal{P}\!\mathit{ic}}
\def\Picardfunctor{\mathrm{Pic}}
\def\Deformationcategory{\mathcal{D}\!\mathit{ef}}
\usepackage{cleveref}
\crefname{originalproposition}{original proposition}{original propositions}
\Crefname{originalproposition}{Original proposition}{Original propositions}
\makeatletter
\newlabel{lemma-hilb-d-separated}{{0B98}{source}{}{sourcepointer-1}{}}
\newlabel{lemma-hilb-d-separated@cref}{{[originalproposition][0][]0B98}{[1][0][]source}{}{}{}}
\makeatother
\crefname{originalproposition}{original proposition}{original propositions}
\Crefname{originalproposition}{Original proposition}{Original propositions}
\makeatletter
\newlabel{lemma-hilb-d-smooth}{{0B9H}{source}{}{sourcepointer-2}{}}
\newlabel{lemma-hilb-d-smooth@cref}{{[originalproposition][0][]0B9H}{[1][0][]source}{}{}{}}
\makeatother
\crefname{originalproposition}{original proposition}{original propositions}
\Crefname{originalproposition}{Original proposition}{Original propositions}
\makeatletter
\newlabel{lemma-universal-object}{{0B9G}{source}{}{sourcepointer-3}{}}
\newlabel{lemma-universal-object@cref}{{[originalproposition][0][]0B9G}{[1][0][]source}{}{}{}}
\makeatother
\crefname{originalproposition}{original proposition}{original propositions}
\Crefname{originalproposition}{Original proposition}{Original propositions}
\makeatletter
\newlabel{morphisms-lemma-image-proper-is-proper}{{03GN}{source}{}{sourcepointer-4}{}}
\newlabel{morphisms-lemma-image-proper-is-proper@cref}{{[originalproposition][0][]03GN}{[1][0][]source}{}{}{}}
\makeatother
\crefname{originalproposition}{original proposition}{original propositions}
\Crefname{originalproposition}{Original proposition}{Original propositions}
\makeatletter
\newlabel{proposition-hilb-d-representable}{{0B9A}{source}{}{sourcepointer-5}{}}
\newlabel{proposition-hilb-d-representable@cref}{{[originalproposition][0][]0B9A}{[1][0][]source}{}{}{}}
\makeatother
\crefname{originalproposition}{original proposition}{original propositions}
\Crefname{originalproposition}{Original proposition}{Original propositions}
\makeatletter
\newlabel{remark-when-proposition-applies}{{0B9B}{source}{}{sourcepointer-6}{}}
\newlabel{remark-when-proposition-applies@cref}{{[originalproposition][0][]0B9B}{[1][0][]source}{}{}{}}
\makeatother
\crefname{originalproposition}{original proposition}{original propositions}
\Crefname{originalproposition}{Original proposition}{Original propositions}
\makeatletter
\newlabel{varieties-lemma-bijection-irreducible-components}{{038F}{source}{}{sourcepointer-7}{}}
\newlabel{varieties-lemma-bijection-irreducible-components@cref}{{[originalproposition][0][]038F}{[1][0][]source}{}{}{}}
\makeatother
\crefname{originalproposition}{original proposition}{original propositions}
\Crefname{originalproposition}{Original proposition}{Original propositions}
\makeatletter
\newlabel{varieties-lemma-dim-1-proper-projective}{{0A26}{source}{}{sourcepointer-8}{}}
\newlabel{varieties-lemma-dim-1-proper-projective@cref}{{[originalproposition][0][]0A26}{[1][0][]source}{}{}{}}
\makeatother
\title{Hilb\textasciicircum{}d of a curve is smooth proper of dimension d with degree formula}
\author{The Stacks Project Authors}
\date{stacks-project:pic.tex:local-snapshot-20261012; GNU FDL 1.2 or later}
\begin{document}
\maketitle
\noindent\textit{Editorial note.} Statement span is lines 848-868 of sources/stacks-pic/pinned-original-source.tex; proof span is lines 870-909, the complete author proof (38 body lines).\par\medskip
\noindent\textit{Editorial note.} Navigation stubs: 8 cited-but-not-included labels are supplied as non-load-bearing pointers with their pinned tags where the snapshot tag index has them.\par\medskip
\noindent\textit{Editorial note.} extra\_packages is empty on purpose: the retained preamble already loads xy with options and a second declaration clashes in the assembler wrapper.\par\medskip
\noindent\textit{Editorial note.} Text is retained verbatim, including the author's own typos and placeholders.\par\medskip
\noindent\textit{Editorial note.} Queue index 11 of queue\_pic.json (pb=1594, stmt 848-868, proof 870-909).\par\medskip
\noindent\textit{Editorial note.} stacks-pic wave-1 (pi01): fresh adjudication accept (2026-10-14); pool per standing queue hypercovering -> pic (SMC reassigned to Zhen per Chief 08:5x); pic.tex 27 labels/23 with proof judged 21 accept / 2 reject / 4 pb=0 env-dup rows excluded; spans cite-free and reference/slogan/history-free so no patchenv needed.\par\medskip
\noindent\textit{Allowed standard prerequisites (editorial).}\begin{itemize}
\item Explicit prerequisite boundary. The cited inputs are prerequisites; the argument in the retained span is not replaced by any cited result.
\end{itemize}\medskip
\ifdefined\hypertarget\hypertarget{sourcepointer-1}{}\fi
\noindent\textit{Source reference.} Original proposition 0B98: sources/stacks-pic/pinned-original-source.tex, source lines 381--385. Cited by the retained statement or proof but not part of the collected spans; supplied only as a navigation pointer. The argument does not depend on its content.\par\medskip
\ifdefined\hypertarget\hypertarget{sourcepointer-2}{}\fi
\noindent\textit{Source reference.} Original proposition 0B9H: sources/stacks-pic/pinned-original-source.tex, source lines 819--824. Cited by the retained statement or proof but not part of the collected spans; supplied only as a navigation pointer. The argument does not depend on its content.\par\medskip
\ifdefined\hypertarget\hypertarget{sourcepointer-3}{}\fi
\noindent\textit{Source reference.} Original proposition 0B9G: sources/stacks-pic/pinned-original-source.tex, source lines 786--792. Cited by the retained statement or proof but not part of the collected spans; supplied only as a navigation pointer. The argument does not depend on its content.\par\medskip
\ifdefined\hypertarget\hypertarget{sourcepointer-4}{}\fi
\noindent\textit{Source reference.} Original proposition 03GN: sources/stacks-pic/pinned-original-source.tex, source lines 909--909. Cited by the retained statement or proof but not part of the collected spans; supplied only as a navigation pointer. The argument does not depend on its content.\par\medskip
\ifdefined\hypertarget\hypertarget{sourcepointer-5}{}\fi
\noindent\textit{Source reference.} Original proposition 0B9A: sources/stacks-pic/pinned-original-source.tex, source lines 516--522. Cited by the retained statement or proof but not part of the collected spans; supplied only as a navigation pointer. The argument does not depend on its content.\par\medskip
\ifdefined\hypertarget\hypertarget{sourcepointer-6}{}\fi
\noindent\textit{Source reference.} Original proposition 0B9B: sources/stacks-pic/pinned-original-source.tex, source lines 589--601. Cited by the retained statement or proof but not part of the collected spans; supplied only as a navigation pointer. The argument does not depend on its content.\par\medskip
\ifdefined\hypertarget\hypertarget{sourcepointer-7}{}\fi
\noindent\textit{Source reference.} Original proposition 038F: sources/stacks-pic/pinned-original-source.tex, source lines 909--909. Cited by the retained statement or proof but not part of the collected spans; supplied only as a navigation pointer. The argument does not depend on its content.\par\medskip
\ifdefined\hypertarget\hypertarget{sourcepointer-8}{}\fi
\noindent\textit{Source reference.} Original proposition 0A26: sources/stacks-pic/pinned-original-source.tex, source lines 909--909. Cited by the retained statement or proof but not part of the collected spans; supplied only as a navigation pointer. The argument does not depend on its content.\par\medskip
\begin{proposition}
\label{proposition-hilb-d}
Let $X$ be a geometrically irreducible smooth proper curve over a field $k$.
\begin{enumerate}
\item The functors $\Hilbfunctor^d_{X/k}$ are representable by smooth
proper varieties $\underline{\Hilbfunctor}^d_{X/k}$ of dimension
$d$ over $k$.
\item For a field extension $k'/k$ the $k'$-rational points
of $\underline{\Hilbfunctor}^d_{X/k}$ are in $1$-to-$1$ bijection
with effective Cartier divisors of degree $d$ on $X_{k'}$.
\item For $d_1, d_2 \geq 0$ there is a morphism
$$
\underline{\Hilbfunctor}^{d_1}_{X/k}
\times_k
\underline{\Hilbfunctor}^{d_2}_{X/k}
\longrightarrow
\underline{\Hilbfunctor}^{d_1 + d_2}_{X/k}
$$
which is finite locally free of degree ${d_1 + d_2 \choose d_1}$.
\end{enumerate}
\end{proposition}

\begin{proof}
The functors $\Hilbfunctor^d_{X/k}$ are representable by
Proposition \ref{proposition-hilb-d-representable}
(see also Remark \ref{remark-when-proposition-applies})
and the fact that $X$ is projective
(Varieties, Lemma \ref{varieties-lemma-dim-1-proper-projective}).
The schemes $\underline{\Hilbfunctor}^d_{X/k}$ are separated
over $k$ by Lemma \ref{lemma-hilb-d-separated}.
The schemes $\underline{\Hilbfunctor}^d_{X/k}$ are smooth
over $k$ by Lemma \ref{lemma-hilb-d-smooth}.
Starting with $X = \underline{\Hilbfunctor}^1_{X/k}$,
the morphisms of Lemma \ref{lemma-universal-object},
and induction we find a morphism
$$
X^d = X \times_k X \times_k \ldots \times_k X \longrightarrow
\underline{\Hilbfunctor}^d_{X/k},\quad
(x_1, \ldots, x_d) \longrightarrow x_1 + \ldots + x_d
$$
which is finite locally free of degree $d!$. Since $X$ is
proper over $k$, so is $X^d$, hence
$\underline{\Hilbfunctor}^d_{X/k}$ is proper over $k$ by
Morphisms, Lemma \ref{morphisms-lemma-image-proper-is-proper}.
Since $X$ is geometrically irreducible over $k$, the product
$X^d$ is irreducible
(Varieties, Lemma \ref{varieties-lemma-bijection-irreducible-components})
hence the image is irreducible (in fact geometrically irreducible).
This proves (1). Part (2) follows from the definitions. Part (3) follows
from the commutative diagram
$$
\xymatrix{
X^{d_1} \times_k X^{d_2} \ar[d] \ar@{=}[r] & X^{d_1 + d_2} \ar[d] \\
\underline{\Hilbfunctor}^{d_1}_{X/k}
\times_k
\underline{\Hilbfunctor}^{d_2}_{X/k}
\ar[r] &
\underline{\Hilbfunctor}^{d_1 + d_2}_{X/k}
}
$$
and multiplicativity of degrees of finite locally free morphisms.
\end{proof}

\end{document}
